% =====================================================================
% Supplementary Information — Nature Climate Change (anonymised)
% =====================================================================
\documentclass[12pt,a4paper]{article}
\usepackage[margin=2.5cm]{geometry}
\usepackage[T1]{fontenc}
\usepackage[utf8]{inputenc}
\usepackage{lmodern}
\usepackage{booktabs}
\usepackage{amsmath}
\usepackage{setspace}
\usepackage[super,sort&compress]{natbib}
\usepackage[hidelinks]{hyperref}
\usepackage{caption}
\captionsetup{font=small,labelfont=bf}
\setlength{\bibsep}{2pt}
\onehalfspacing
\newcommand{\dd}{\ensuremath{\,^{\circ}\mathrm{C}}}

\begin{document}
\begin{center}
{\LARGE\bfseries Supplementary Information}\\[6pt]
{\large Night-time moist heat driven by humidity, not temperature alone, is crossing
physiologically anchored survival thresholds}\\[4pt]
{\small [Author list withheld for double-blind peer review]}
\end{center}

\section*{Supplementary Note 1: Data processing and quality control}
\subsection*{Station screening}
GHCNh version 1.1.0 provides sub-hourly records for 5,021 station archives; 176 stations satisfy
the a-priori screen of at least 20 station-years with usable records within 1976--2023. Of 4,763
downloaded station-years, 4,164 (87\%) pass the quality gate: dry-bulb and dew-point present,
physical range checks ($-60<T<60$\,\dd{}; $-80<T_\mathrm{d}<40$\,\dd{}), and within-record
plausibility. Wet-bulb temperature is recomputed from dry-bulb plus dew point rather than taken
from native fields because native wet-bulb availability is only 40.6\% and drifts strongly by
decade (for example, one Andean station reports zero usable native wet-bulb values in 1976 versus
full availability after re-computation). Three wet-bulb implementations (Stull; Davies--Jones;
dew-point inversion) agree to $0.05\,\dd$ in the decision-relevant range.

\subsection*{Night definition}
Night is the set of hours with solar elevation $<-0.833^\circ$ computed from station latitude,
longitude and UTC timestamp. Clock-hour definitions (for example fixed 22:00--06:00 local)
underestimate $\mathrm{NWH}_{21.9}$ by 27.7--31.5\% and dominate the implementation uncertainty
budget; the astronomical definition is therefore part of the metric specification, and the ratio
of ratio-type indicators between definitions is 1.00.

\subsection*{Population weighting}
Station populations are computed from GeoNames city points: full weight within 50\,km, exponential
decay to 300\,km. Total network coverage is 9.01 billion person-equivalents (with overlap);
population weighting raises the headline trend from $+42.84$ to $+52.13$\,h per decade.

\section*{Supplementary Note 2: Negative and null diagnostics}
\subsection*{WHO Mortality Database annual panels}
Annual all-cause mortality panels for 12 countries (238 country-years) show no wet-tropics
humid-heat effect in the ICD-10 era ($z=+0.08$); the apparent effect in the full-period panel is
generated by early registration-coverage changes. Because annual aggregation also dilutes the
European weekly signal by $\sim$93\% (z from $+34.4$ to $+1.1$), the WHO panels are excluded from
inference and retained only as a negative diagnostic of aggregation plus registration artefacts.

\subsection*{First-crossing hazard model (retracted)}
A time-to-first-crossing model across 64 European regions returned HR $=0.863$ $[0.603,1.236]$,
which we retract: the at-risk set cools by $-2.50\,\dd$ (median background night-time
$T_\mathrm{w}$ falls from 11.75 to 9.26\,\dd) because the most humid regions cross first and leave
the risk set. This is frailty depletion, not protection; the analysis is retained in the frozen
artefact set solely for transparency (Extended Data Table~1, row R1).

\subsection*{Singapore NWH21.9 exposure--response (rejected artefact)}
With a median monthly exceedance fraction of 1.000, $\mathrm{NWH}_{21.9}$ in Singapore is
collinear with night length; the model returns $z=+8.61$, which we reject as a saturation
artefact. All wet-tropics inference uses $\mathrm{NWH}_{25.0}$/$\mathrm{NWH}_{25.8}$ or
within-city quantiles.

\section*{Supplementary Note 3: Register of retracted, excluded and corrected analyses}
During 27 rounds of internal adversarial review, the following were retracted, excluded or
corrected and are never presented as positive results (see also Extended Data Table~1, rows
R1--R6):
\begin{enumerate}
\item \textbf{R1 (retracted):} first-crossing hazard ratio --- frailty depletion (Note 2).
\item \textbf{R2 (excluded):} WHO-style annual mortality panel --- registration artefact plus
annual aggregation (Note 2).
\item \textbf{R3 (rejected):} Singapore $\mathrm{NWH}_{21.9}$ exposure--response --- saturation
artefact (Note 2).
\item \textbf{R4 (corrected):} a formatting bug printed $I^2$ as 1\% instead of 78.2\%; all
published values quote the frozen JSON.
\item \textbf{R5 (corrected):} a two-cluster quantile model reported a cluster-robust
$z=+12.97$; independent HC0 intervals are reported instead (CI includes 1).
\item \textbf{R6 (corrected):} an earlier draft stated that exposure extent grows faster than
intensity; the opposite holds (intensity $+47\%$ vs extent $+26\%$), and the $1.78\times$
overstatement factor follows because risk tracks extent.
\end{enumerate}
All 14 core scripts re-run against frozen baselines reproduce byte-identical results (SHA-256,
25\,min\,36\,s total), and every displayed number in the Article is registered in a manifest
tracing it to a hash-verified artefact.

\section*{Supplementary Note 4: Robustness summary}
\begin{table}[htbp]
\centering\small
\caption{Key robustness results (all values from frozen artefacts).}
\begin{tabular}{p{7.2cm}p{7.6cm}}
\toprule
Check & Result \\
\midrule
Analysis window (8 windows, 1976--2023) & $+27.69$ to $+58.62$\,h per decade, all positive \\
Aggregation calibre & mean $+42.84$ / population-weighted $+52.13$ / median $+18.27$\,h per decade \\
Homogenisation & trend ratio 1.00 after neighbour-reference homogenisation \\
Urban heat island & most-rural quartile remains significant ($p=6.8\times10^{-4}$) \\
Wet-bulb implementation & Stull / Davies--Jones / inversion agree to $0.05\,\dd$ \\
External confirmation & HadEX3 TN90p vs night-recovery failure: $r=0.402$ ($p=0.0069$), 92\% sign agreement; GISTEMP 91\% sign agreement ($r=0.490$) \\
Effective sample size & $N_\mathrm{eff}=67$ ($\bar r=-0.015$); cluster bootstrap CI for median trend $[+3.4,+57.0]$\,h per decade \\
Lag-1 mortality displacement & RR $=0.9979$ (no harvesting) \\
Exposure--response gradient & positive unrestricted; negative restricted to informative regions (ED Fig.~5) \\
Conservatism & raw station trends $\approx$ 0.46--0.85$\times$ homogenised gridded products \\
\bottomrule
\end{tabular}
\end{table}

\section*{Supplementary references}
\begin{enumerate}
\item Raymond, C., Matthews, T. \& Horton, R. M. The emergence of heat and humidity too severe for
human tolerance. \textit{Sci. Adv.} \textbf{6}, eaaw1838 (2020).
\item Vanos, J. \textit{et al.} A physiological approach for assessing human survivability and
liveability to heat in a changing climate. \textit{Nat. Commun.} \textbf{14}, 7653 (2023).
\item Ballester, J. \textit{et al.} The effect of temporal data aggregation to assess the impact of
changing temperatures in Europe. \textit{Lancet Reg. Health Eur.} \textbf{36}, 100779 (2024).
\item World Meteorological Organization \& World Health Organization. \textit{Heatwaves and
Health: Guidance on Warning-System Development} (WMO-No. 1142, 2015).
\end{enumerate}

\end{document}
